% !TEX TS-program = xelatex
% !TEX TS-program = xelatex

\documentclass[12pt,a4paper]{article}

\IfFileExists{src/libs/body.tex}{% --- body.tex ---
\usepackage{pdfpages}
\usepackage{fontspec}
% Prefer Times New Roman when available; fall back in container environments.
\IfFontExistsTF{Times New Roman}{
  \setmainfont{Times New Roman}
}{
  \setmainfont{Liberation Serif}
}
\usepackage[vietnamese]{babel}

% Page layout and spacing
\usepackage[left=3cm,right=2cm,top=2cm,bottom=2cm]{geometry}
\usepackage{titlesec}
\usepackage{setspace}

% Silence noisy warnings
\usepackage{silence}
\WarningFilter{transparent}{Loading aborted}
\WarningFilter{hyperref}{The anchor of a bookmark and its parent's}

% Core packages
\usepackage{float}
\usepackage{svg}
\usepackage{graphicx}
\graphicspath{{assets/}{../assets/}{../../assets/}}
\svgpath{{assets/}{../assets/}{../../assets/}}
\usepackage{enumitem}
\usepackage{array}
\usepackage{tabularx}
\usepackage{multirow}
\usepackage[table]{xcolor}
\usepackage{ulem}
\usepackage{xurl}

% ====== THÊM 2 GÓI NÀY ĐỂ FIX LỖI ADJUSTBOX VÀ FOREST ======
\usepackage{adjustbox}
\usepackage[edges]{forest}
% ============================================================

% TikZ
\usepackage{tikz}

% ====== BỔ SUNG FIT, BACKGROUNDS, TREES VÀO DANH SÁCH ======
\usetikzlibrary{calc, arrows.meta, shapes.geometric, positioning, fit, backgrounds, trees}

% Table tuning
\renewcommand{\tabularxcolumn}[1]{m{#1}}
\hbadness=10000
\hfuzz=10000pt

% Shared commands
\newcommand{\subsecheading}[1]{%
    \par\vspace{2pt}%
    \hspace{1.5em}\textbf{#1}%
    \par\vspace{2pt}%
}
\newcommand{\introtext}[1]{%
    \par\hspace{3.0em}#1\par%
}
\newcommand{\StandaloneDocumentSetup}[2]{%
  \pagenumbering{arabic}%
  \ApplyStandardFooter{#1}{#2}%
}
\newcommand{\StandaloneSectionSetup}[2]{%
  \ApplyStandardFooter{#1}{#2}%
}

% Math
\usepackage{amsmath}

% Hyperref
\usepackage{hyperref}
\hypersetup{
    colorlinks=true,
    linkcolor=black,
    filecolor=magenta,
    urlcolor=blue,
}}{% --- body.tex ---
\usepackage{pdfpages}
\usepackage{fontspec}
% Prefer Times New Roman when available; fall back in container environments.
\IfFontExistsTF{Times New Roman}{
  \setmainfont{Times New Roman}
}{
  \setmainfont{Liberation Serif}
}
\usepackage[vietnamese]{babel}

% Page layout and spacing
\usepackage[left=3cm,right=2cm,top=2cm,bottom=2cm]{geometry}
\usepackage{titlesec}
\usepackage{setspace}

% Silence noisy warnings
\usepackage{silence}
\WarningFilter{transparent}{Loading aborted}
\WarningFilter{hyperref}{The anchor of a bookmark and its parent's}

% Core packages
\usepackage{float}
\usepackage{svg}
\usepackage{graphicx}
\graphicspath{{assets/}{../assets/}{../../assets/}}
\svgpath{{assets/}{../assets/}{../../assets/}}
\usepackage{enumitem}
\usepackage{array}
\usepackage{tabularx}
\usepackage{multirow}
\usepackage[table]{xcolor}
\usepackage{ulem}
\usepackage{xurl}

% ====== THÊM 2 GÓI NÀY ĐỂ FIX LỖI ADJUSTBOX VÀ FOREST ======
\usepackage{adjustbox}
\usepackage[edges]{forest}
% ============================================================

% TikZ
\usepackage{tikz}

% ====== BỔ SUNG FIT, BACKGROUNDS, TREES VÀO DANH SÁCH ======
\usetikzlibrary{calc, arrows.meta, shapes.geometric, positioning, fit, backgrounds, trees}

% Table tuning
\renewcommand{\tabularxcolumn}[1]{m{#1}}
\hbadness=10000
\hfuzz=10000pt

% Shared commands
\newcommand{\subsecheading}[1]{%
    \par\vspace{2pt}%
    \hspace{1.5em}\textbf{#1}%
    \par\vspace{2pt}%
}
\newcommand{\introtext}[1]{%
    \par\hspace{3.0em}#1\par%
}
\newcommand{\StandaloneDocumentSetup}[2]{%
  \pagenumbering{arabic}%
  \ApplyStandardFooter{#1}{#2}%
}
\newcommand{\StandaloneSectionSetup}[2]{%
  \ApplyStandardFooter{#1}{#2}%
}

% Math
\usepackage{amsmath}

% Hyperref
\usepackage{hyperref}
\hypersetup{
    colorlinks=true,
    linkcolor=black,
    filecolor=magenta,
    urlcolor=blue,
}}
\IfFileExists{src/libs/header.tex}{% --- header.tex ---
\usepackage{fancyhdr}
\setlength{\headheight}{15.5pt} % Thêm dòng này để fix cảnh báo chiều cao
%\setlength{\headheight}{14.5pt}
\addtolength{\topmargin}{-2.5pt}

\pagestyle{fancy}
\fancyhf{}
\renewcommand{\headrulewidth}{0pt}
\renewcommand{\footrulewidth}{0pt}

\newcommand{\ApplyStandardHeader}{%
  \fancyhead[L]{%
    \begin{minipage}[t]{0.795\linewidth}%
      UIT \textbar{} SS004.F21 \textbar{} Nhóm 03 \textbar{} Đồ án cuối kỳ%
    \end{minipage}%
    \vrule%
    \begin{minipage}[t]{0.185\linewidth}%
      \centering cTetris%
    \end{minipage}%
  }%
}

\ApplyStandardHeader
}{% --- header.tex ---
\usepackage{fancyhdr}
\setlength{\headheight}{15.5pt} % Thêm dòng này để fix cảnh báo chiều cao
%\setlength{\headheight}{14.5pt}
\addtolength{\topmargin}{-2.5pt}

\pagestyle{fancy}
\fancyhf{}
\renewcommand{\headrulewidth}{0pt}
\renewcommand{\footrulewidth}{0pt}

\newcommand{\ApplyStandardHeader}{%
  \fancyhead[L]{%
    \begin{minipage}[t]{0.795\linewidth}%
      UIT \textbar{} SS004.F21 \textbar{} Nhóm 03 \textbar{} Đồ án cuối kỳ%
    \end{minipage}%
    \vrule%
    \begin{minipage}[t]{0.185\linewidth}%
      \centering cTetris%
    \end{minipage}%
  }%
}

\ApplyStandardHeader
}
\IfFileExists{src/libs/footer.tex}{% --- footer.tex ---
\usepackage{lastpage}

\newcommand{\ApplyStandardFooter}[2]{%
  \fancyfoot[L]{%
    \begin{minipage}[t]{0.795\linewidth}%
      Document ID: #1%
    \end{minipage}%
    \vrule%
    \begin{minipage}[t]{0.185\linewidth}%
      \centering\thepage/\pageref{#2}%
    \end{minipage}%
  }%
  \fancypagestyle{plain}{%
    \fancyhf{}%
    \renewcommand{\headrulewidth}{0pt}%
    \renewcommand{\footrulewidth}{0pt}%
    \ApplyStandardHeader%
    \fancyfoot[L]{%
      \begin{minipage}[t]{0.795\linewidth}%
        Document ID: #1%
      \end{minipage}%
      \vrule%
      \begin{minipage}[t]{0.185\linewidth}%
        \centering\thepage/\pageref{#2}%
      \end{minipage}%
    }%
  }%
  \pagestyle{fancy}%
}
}{% --- footer.tex ---
\usepackage{lastpage}

\newcommand{\ApplyStandardFooter}[2]{%
  \fancyfoot[L]{%
    \begin{minipage}[t]{0.795\linewidth}%
      Document ID: #1%
    \end{minipage}%
    \vrule%
    \begin{minipage}[t]{0.185\linewidth}%
      \centering\thepage/\pageref{#2}%
    \end{minipage}%
  }%
  \fancypagestyle{plain}{%
    \fancyhf{}%
    \renewcommand{\headrulewidth}{0pt}%
    \renewcommand{\footrulewidth}{0pt}%
    \ApplyStandardHeader%
    \fancyfoot[L]{%
      \begin{minipage}[t]{0.795\linewidth}%
        Document ID: #1%
      \end{minipage}%
      \vrule%
      \begin{minipage}[t]{0.185\linewidth}%
        \centering\thepage/\pageref{#2}%
      \end{minipage}%
    }%
  }%
  \pagestyle{fancy}%
}
}
\usepackage{docmute}
\hypersetup{pdftitle={Báo cáo Kỹ năng nghề nghiệp - chiến lược kiểm thử}}

\begin{document}
\providecommand{\StandaloneSectionSetup}[2]{\ApplyStandardFooter{#1}{#2}}
\StandaloneSectionSetup{11-testStrategy.tex}{LastPageTestStrategy}
\onehalfspacing
\section{Chiến Lược Kiểm Thử}

% ── 1. GIỚI THIỆU ────────────────────────────────────────────────────
\subsection{Giới Thiệu}

\subsubsection{Mục Đích Tài Liệu}
\introtext{Tài liệu này định nghĩa chiến lược kiểm thử tổng thể cho dự án cTetris – game xếp hình viết bằng C++ chạy trên console. Chiến lược được xây dựng phù hợp với môi trường phát triển nhóm đa lập trình viên, bao gồm các phương pháp kiểm thử tiêu chuẩn trong ngành game development nhằm đảm bảo chất lượng sản phẩm từ giai đoạn coding đến release.}

\subsubsection{Phạm Vi Áp Dụng}
\introtext{Chiến lược áp dụng cho toàn bộ ba module cấu thành hệ thống:}
\begin{itemize}[leftmargin=3em]
  \item \textbf{gameCore} – Logic gameplay lõi: xử lý Tetrimino, va chạm, tính điểm, quản lý trạng thái
  \item \textbf{gameConsole} – Giao diện console: xử lý đầu vào, rendering, game loop
  \item \textbf{gameStory} – Nội dung narrative: sự kiện cốt truyện, giao diện giới thiệu
\end{itemize}

\subsubsection{Tài Liệu Tham Chiếu}
\begin{itemize}[leftmargin=3em]
  \item \texttt{app/appTree.md} – Cấu trúc module và dependency của dự án
  \item GDC (Game Developers Conference) QA Testing Guidelines
  \item IGDA Quality Assurance Best Practices
  \item Google Test \& Catch2 Framework Documentation
\end{itemize}

% ── 2. MỤC TIÊU KIỂM THỬ ────────────────────────────────────────────
\subsection{Mục Tiêu Kiểm Thử}

\introtext{Chiến lược kiểm thử hướng tới sáu mục tiêu cốt lõi:}

\begin{table}[H]
\centering
\rowcolors{2}{gray!15}{white}
\begin{tabularx}{\linewidth}{|c|X|X|}
\hline
\rowcolor{gray!35}
\textbf{\#} & \textbf{Mục Tiêu} & \textbf{Phương Pháp Kiểm Thử} \\ \hline
MT-01 & Đảm bảo tính đúng đắn của luật chơi Tetris (Tetrimino, xoay, va chạm, xóa dòng) & Unit Test, System Test \\ \hline
MT-02 & Đảm bảo tương tác thời gian thực không có lag hay delay ảnh hưởng trải nghiệm & Performance Test, Usability Test \\ \hline
MT-03 & Giảm thiểu lỗi tích hợp trong môi trường phát triển nhóm đa lập trình viên & Integration Test, CI/CD \\ \hline
MT-04 & Duy trì tính nhất quán của giao diện console và trạng thái game xuyên suốt phiên chơi & System Test, Regression Test \\ \hline
MT-05 & Phát hiện sớm lỗi hồi quy khi thêm tính năng mới (tốc độ, chế độ chơi) & Automated Regression Test \\ \hline
MT-06 & Đảm bảo hiệu năng ổn định và khả năng mở rộng trên nhiều cấu hình & Performance Test, Load Test \\ \hline
\end{tabularx}
\end{table}

% ── 3. PHẠM VI KIỂM THỬ ─────────────────────────────────────────────
\subsection{Phạm Vi Kiểm Thử}

\subsubsection{Module gameCore \textnormal{(Mức Ưu Tiên: Cao)}}
\introtext{gameCore là module lõi xử lý toàn bộ logic gameplay. Phạm vi kiểm thử bao gồm:}
\begin{itemize}[leftmargin=3em]
  \item Tạo và quản lý 7 loại Tetrimino (I, O, T, S, Z, J, L) theo bag randomizer
  \item Di chuyển và xoay theo chuẩn Super Rotation System (SRS), bao gồm wall kick
  \item Phát hiện va chạm chính xác: đáy, thành bên, khối đã đặt; không cho phép overlap
  \item Xóa dòng và tính điểm theo công thức chuẩn: Single / Double / Triple / Tetris
  \item Quản lý vòng đời game: playing $\rightarrow$ paused $\rightarrow$ game over $\rightarrow$ restart
  \item Tiến trình level và tốc độ rơi tăng dần sau mỗi 10 dòng được xóa
  \item Điều kiện biên: xoay ở góc, game over khi vùng spawn bị chặn
\end{itemize}

\subsubsection{Module gameConsole \textnormal{(Mức Ưu Tiên: Trung Bình)}}
\introtext{gameConsole quản lý giao diện và tương tác người dùng. Phạm vi kiểm thử:}
\begin{itemize}[leftmargin=3em]
  \item Xử lý đầu vào bàn phím chuẩn: di chuyển, xoay, Space (hard drop), P (pause)
  \item Hiển thị ASCII art board 10$\times$20, next piece, score, level, lines cleared chính xác
  \item Game loop ổn định (target 60 FPS), không flicker hay lag
  \item Quản lý trạng thái UI: menu pause, game over screen, high score, transitions
  \item Tương thích terminal: không wrap text, không color/encoding issues
  \item Xử lý hold keys liên tục, không miss input trong fast gameplay
\end{itemize}

\subsubsection{Module gameStory \textnormal{(Mức Ưu Tiên: Thấp)}}
\introtext{gameStory quản lý nội dung narrative và onboarding. Phạm vi kiểm thử:}
\begin{itemize}[leftmargin=3em]
  \item Kích hoạt event dựa trên tiến trình game (unlock level, hints, cutscenes)
  \item Hiển thị text và transitions không gián đoạn gameplay đang chạy
  \item Đảm bảo story không interfere với logic gameCore
  \item Scalability: sẵn sàng mở rộng khi content được bổ sung
\end{itemize}

\subsubsection{Kiểm Thử Tích Hợp}
\begin{itemize}[leftmargin=3em]
  \item gameCore $\leftrightarrow$ gameConsole: input $\rightarrow$ logic update $\rightarrow$ render phải đồng bộ chính xác
  \item gameCore $\leftrightarrow$ gameStory: event level up trigger story element không gây lag
  \item End-to-end scenario: spawn $\rightarrow$ play $\rightarrow$ pause/resume $\rightarrow$ scoring $\rightarrow$ level up $\rightarrow$ game over
  \item Stress testing: phiên dài phát hiện memory leaks và performance degradation
\end{itemize}

% ── 4. CHIẾN LƯỢC KIỂM THỬ THEO LỚP ─────────────────────────────────
\subsection{Chiến Lược Kiểm Thử Theo Lớp}

\introtext{Chiến lược áp dụng mô hình Test Pyramid, từ unit testing đến system testing, kết hợp kiểm thử tự động và thủ công theo tỷ lệ phù hợp với game development.}

\begin{table}[H]
\centering
\rowcolors{2}{gray!15}{white}
\begin{tabularx}{\linewidth}{|l|X|X|X|}
\hline
\rowcolor{gray!35}
\textbf{Tầng} & \textbf{Phương Pháp / Đối Tượng} & \textbf{Công Cụ} & \textbf{Tự Động} \\ \hline
Unit        & Kiểm thử đơn vị từng hàm/component – gameCore, Console, Story             & Google Test / Catch2     & 100\%   \\ \hline
Integration & Kiểm thử tương tác giữa các module – Core$\leftrightarrow$Console, Core$\leftrightarrow$Story & Stubs/Mocks + Real Runs  & 70\%    \\ \hline
System      & End-to-end flow, stress, compatibility – Toàn bộ hệ thống                 & Manual + Script          & 40\%    \\ \hline
Regression  & Chạy sau mỗi commit, bảo vệ core – Core mechanics, state                  & CI/CD Pipeline           & 100\%   \\ \hline
Performance & FPS, latency, memory, CPU – Game loop, input                              & Benchmark tools          & Partial \\ \hline
Usability   & Playtesting, balance, controls – Toàn bộ game                             & Manual / Người chơi thật & 0\%     \\ \hline
\end{tabularx}
\end{table}

\subsubsection{Unit Testing}
\begin{itemize}[leftmargin=3em]
  \item \textbf{gameCore:} \texttt{rotate()}, \texttt{detectCollision()}, \texttt{clearLine()}, \texttt{calculateScore()}, \texttt{spawnPiece()}
  \item \textbf{gameConsole:} \texttt{parseInput()}, \texttt{renderBoard()}, \texttt{renderScore()}, \texttt{updateUI()}
  \item \textbf{gameStory:} \texttt{triggerEvent()}, \texttt{loadDialogue()}, \texttt{checkProgress()}
  \item Nguyên tắc: mỗi test độc lập, assert rõ ràng, không chỉ kiểm tra không crash
\end{itemize}

\subsubsection{Integration Testing}
\begin{itemize}[leftmargin=3em]
  \item gameCore + gameConsole: input từ console cập nhật logic đúng, display phản ánh state chính xác
  \item gameCore + gameStory: event trigger không làm gián đoạn gameplay
  \item Phương pháp: stubs/mocks cho isolated testing, real runs cho full integration
\end{itemize}

\subsubsection{System Testing}
\begin{itemize}[leftmargin=3em]
  \item Playthrough scenario hoàn chỉnh từ spawn đến game over
  \item Stress testing: session 30+ phút kiểm tra stability và memory
  \item Compatibility: các terminal emulator và OS versions khác nhau
\end{itemize}

\subsubsection{Regression Testing}
\begin{itemize}[leftmargin=3em]
  \item Automated suite chạy tự động sau mỗi commit via CI/CD
  \item Bao phủ core mechanics: Tetrimino spawning, scoring, state transitions
  \item Kết quả phải 100\% pass trước khi merge vào main branch
\end{itemize}

\subsubsection{Performance Testing}
\begin{itemize}[leftmargin=3em]
  \item Frame rate: đảm bảo $\geq$ 60 FPS liên tục, không drop trong fast play
  \item Input latency: đo response time từ key press đến game action (target $<$ 16ms)
  \item Memory: monitor trong long sessions, không memory leak
  \item Benchmarking: so sánh chỉ số giữa các build version
\end{itemize}

\subsubsection{Usability Testing (Playtesting)}
\begin{itemize}[leftmargin=3em]
  \item Controls trực quan và dễ học, menu dễ navigate
  \item Difficulty curve cân bằng, scoring cảm thấy công bằng
  \item Thời gian phản hồi phím trực tiếp, không có delay cảm nhận được
\end{itemize}

% ── 5. CHIẾN LƯỢC CHO MÔI TRƯỜNG ĐA LẬP TRÌNH VIÊN ──────────────────
\subsection{Chiến Lược Cho Môi Trường Đa Lập Trình Viên}

\subsubsection{Phân Chia Trách Nhiệm Kiểm Thử}

\begin{table}[H]
\centering
\rowcolors{2}{gray!15}{white}
\begin{tabularx}{\linewidth}{|l|l|X|X|}
\hline
\rowcolor{gray!35}
\textbf{Module} & \textbf{Developer Phụ Trách} & \textbf{Phạm Vi Kiểm Thử} & \textbf{Giao Diện Kiểm Tra} \\ \hline
gameCore    & Developer Core    & Logic, collision, scoring, state & API hàm được export \\ \hline
gameConsole & Developer Console & UI, input, rendering, game loop  & Key events, output buffer \\ \hline
gameStory   & Developer Story   & Events, dialogue, content loading & Event callbacks, JSON \\ \hline
Integration & Toàn nhóm (Tech Lead) & Cross-module communication    & End-to-end test cases \\ \hline
\end{tabularx}
\end{table}

\subsubsection{Giao Thức Kiểm Thử Chung}
\begin{itemize}[leftmargin=3em]
  \item Quy ước đặt tên test case: \texttt{[MODULE]\_TC\_[ID]\_[TenNganGon]}
  \item Ví dụ: \texttt{CORE\_TC\_001\_RotateBlock}, \texttt{CON\_TC\_005\_PauseResume}
  \item Mỗi test case bắt buộc ghi: Preconditions, Steps, Expected Output, Actual Output, Pass/Fail
  \item Ranh giới module được định nghĩa rõ ràng qua API documentation trước khi code
  \item Mọi commit feature phải kèm test case đã chạy và kết quả
\end{itemize}

\subsubsection{Quy Trình Pull Request}
\begin{itemize}[leftmargin=3em]
  \item Developer chạy unit test của module mình, xác nhận 100\% pass
  \item PR mô tả rõ test case đã thực hiện, bao gồm edge case
  \item Bắt buộc chạy integration test trước khi merge
  \item CI/CD tự động chạy regression suite – PR fail CI không được merge
  \item Reviewer xác minh test scope và kết quả trước khi approve
\end{itemize}

% ── 6. KIỂM THỬ TỰ ĐỘNG ──────────────────────────────────────────────
\subsection{Kiểm Thử Tự Động}

\subsubsection{Framework và Công Cụ}

\begin{table}[H]
\centering
\rowcolors{2}{gray!15}{white}
\begin{tabularx}{\linewidth}{|l|X|l|X|}
\hline
\rowcolor{gray!35}
\textbf{Công Cụ} & \textbf{Mục Đích} & \textbf{Phạm Vi} & \textbf{Ghi Chú} \\ \hline
Google Test      & Unit \& Integration Testing  & Toàn bộ modules     & Recommended primary \\ \hline
Catch2           & Unit Testing (alternative)   & gameCore functions  & Syntax đơn giản hơn \\ \hline
CI/CD Pipeline   & Tự động regression test      & Mỗi commit/PR       & GitHub Actions / Jenkins \\ \hline
gcov / llvm-cov  & Code coverage analysis       & Unit test report    & Mục tiêu $\geq$ 80\% core \\ \hline
\end{tabularx}
\end{table}

\subsubsection{Nguyên Tắc Viết Test Tự Động}
\begin{itemize}[leftmargin=3em]
  \item Tách biệt hoàn toàn logic (gameCore) khỏi display (gameConsole): test core không cần terminal
  \item Sử dụng fixed board states và mock data để test deterministic
  \item Mỗi test case độc lập, kết quả không phụ thuộc thứ tự thực thi
  \item Assert phải có ý nghĩa, kiểm tra cả positive case và negative case
  \item Test phải chạy được trong $<$ 10 giây (không block CI pipeline)
\end{itemize}

\subsubsection{Mục Tiêu Code Coverage}
\begin{itemize}[leftmargin=3em]
  \item gameCore (collision, line clear, scoring): tối thiểu 85\% line coverage
  \item gameConsole (input parsing, render logic): tối thiểu 70\% line coverage
  \item gameStory (event trigger logic): tối thiểu 60\% line coverage
\end{itemize}

% ── 7. KIỂM THỬ THỦ CÔNG ─────────────────────────────────────────────
\subsection{Kiểm Thử Thủ Công}

\subsubsection{Phạm Vi Kiểm Thử Thủ Công}
\begin{itemize}[leftmargin=3em]
  \item Trải nghiệm gameplay tổng thể: cảm giác, nhịp độ, fun factor
  \item Giao diện console visual trên các terminal thực tế
  \item Edge case khó tái tạo tự động: nhập phím liên tục, nhiều phím cùng lúc
  \item Âm thanh và hiệu ứng animation (nếu có trong gameStory)
\end{itemize}

\subsubsection{Cấu Trúc Test Case Thủ Công}

\introtext{Mỗi test case thủ công phải có đủ các thông tin sau:}

\begin{table}[H]
\centering
\rowcolors{2}{gray!15}{white}
\begin{tabularx}{\linewidth}{|l|X|}
\hline
\rowcolor{gray!35}
\textbf{Trường} & \textbf{Nội Dung} \\ \hline
Mã Test Case    & \texttt{[MODULE]\_TC\_[ID]\_[TenNganGon]} \\ \hline
Tiêu đề         & Mô tả ngắn gọn chức năng đang test \\ \hline
Preconditions   & Trạng thái hệ thống trước khi chạy test \\ \hline
Steps           & Các bước thực hiện theo thứ tự (đánh số) \\ \hline
Expected Result & Kết quả mong đợi, có thể đo được \\ \hline
Actual Result   & Kết quả thực tế sau khi thực hiện \\ \hline
Pass / Fail     & Kết quả đánh giá và ghi chú nếu fail \\ \hline
\end{tabularx}
\end{table}

\subsubsection{Exploratory Testing}
\begin{itemize}[leftmargin=3em]
  \item Thực hiện bởi thành viên không phải tác giả module để tránh blind spot
  \item Kịch bản: nhập phím bất thường, đổi tốc độ mid-game, session rất dài
  \item Ghi nhận phát hiện theo template bug report chuẩn của nhóm
\end{itemize}

% ── 8. ĐÁNH GIÁ RỦI RO ───────────────────────────────────────────────
\subsection{Đánh Giá Rủi Ro và Ưu Tiên Kiểm Thử}

\begin{table}[H]
\centering
\rowcolors{2}{gray!15}{white}
\begin{tabularx}{\linewidth}{|X|l|l|l|l|}
\hline
\rowcolor{gray!35}
\textbf{Rủi Ro} & \textbf{Module} & \textbf{Mức Độ} & \textbf{P} & \textbf{Kiểu Test} \\ \hline
Sai logic phát hiện va chạm hoặc xếp khối     & gameCore         & Cao         & P1 & Unit        \\ \hline
Tính điểm sai công thức chuẩn Tetris          & gameCore         & Cao         & P1 & Unit        \\ \hline
Game over không đúng điều kiện hoặc restart lỗi & gameCore       & Cao         & P1 & System      \\ \hline
Lỗi tích hợp: state board hiển thị sai so với logic & Core$\leftrightarrow$Console & Cao & P1 & Integration \\ \hline
Điểm số và UI không đồng bộ với state thực    & gameConsole      & Trung bình  & P2 & Integration \\ \hline
Event story kích hoạt sai thời điểm           & gameStory        & Trung bình  & P2 & Integration \\ \hline
Frame drop trong fast gameplay                & gameConsole      & Trung bình  & P2 & Performance \\ \hline
Thẩm mỹ ký tự viền ASCII art bị lỗi          & gameConsole      & Thấp        & P3 & Manual      \\ \hline
Chức năng mở rộng story ảnh hưởng gameplay cơ bản & gameStory   & Thấp        & P3 & System      \\ \hline
\end{tabularx}
\end{table}

% ── 9. QUẢN LÝ TÀI LIỆU VÀ BÁO CÁO ──────────────────────────────────
\subsection{Quản Lý Tài Liệu và Báo Cáo}

\subsubsection{Tài Liệu Test Case}
\introtext{Tài liệu test case được lưu trữ tập trung (repository \texttt{/docs/test-cases/}), cập nhật bắt buộc khi có thay đổi logic. Cấu trúc thư mục theo module: \texttt{/core/}, \texttt{/console/}, \texttt{/story/}, \texttt{/integration/}.}

\subsubsection{Báo Cáo Lỗi (Bug Report)}
\begin{itemize}[leftmargin=3em]
  \item Mô tả hành vi sai (thực tế vs mong đợi)
  \item Bước tái tạo lỗi (Steps to Reproduce) – tối thiểu 3 bước chi tiết
  \item Module liên quan và mức độ ưu tiên: Critical / High / Medium / Low
  \item Test case đã chạy phát hiện lỗi
  \item Screenshot hoặc log terminal nếu có
\end{itemize}

% ── 10. LÝ DO LỰA CHỌN CHIẾN LƯỢC ───────────────────────────────────
\subsection{Lý Do Lựa Chọn Chiến Lược}

\introtext{Chiến lược hiện tại được lựa chọn dựa trên các căn cứ sau:}
\begin{itemize}[leftmargin=3em]
  \item \textbf{Phù hợp kiến trúc module:} kiểm thử từng module riêng biệt giúp phân chia công việc rõ ràng và phát hiện lỗi sớm
  \item \textbf{Hỗ trợ phát triển song song:} mỗi developer tự kiểm thử module của mình, giảm xung đột, tăng hiệu quả phối hợp qua integration test
  \item \textbf{Cân bằng tự động và thủ công:} tự động cho logic lõi (deterministic), thủ công cho UX và visual (cần phán đoán con người)
  \item \textbf{Tập trung vào rủi ro cao:} ưu tiên gameCore và integration, nơi lỗi gây ảnh hưởng lớn nhất đến gameplay
  \item \textbf{Dễ mở rộng:} có thể bổ sung fuzz testing, model-based testing khi dự án lớn hơn
  \item \textbf{Thực tiễn:} dựa trên best practices của ngành game development, đã được kiểm chứng qua các dự án tương tự
  \item \textbf{Không quá phức tạp:} tránh model-based testing làm chính vì đòi hỏi chuyên môn cao không phù hợp scope dự án
\end{itemize}

% ── 11. KẾT LUẬN ─────────────────────────────────────────────────────
\subsection{Kết Luận}

\introtext{Chiến lược kiểm thử này được thiết kế phù hợp với đặc thù của game puzzle Tetris – nơi logic chính xác và trải nghiệm người dùng quyết định chất lượng sản phẩm. Bằng cách áp dụng mô hình kiểm thử phân lớp kết hợp với quy trình CI/CD và phân chia trách nhiệm module rõ ràng, chiến lược đảm bảo:}
\begin{itemize}[leftmargin=3em]
  \item Logic Tetris (Tetrimino handling, scoring, level progression) được kiểm thử toàn diện với coverage cao
  \item Giao diện console đáp ứng tiêu chuẩn usability của game development
  \item Tích hợp module mượt mà, tránh bugs ảnh hưởng đến fun factor
  \item Hồi quy tự động ngăn regression khi codebase phát triển
  \item Hiệu năng ổn định, phù hợp với yêu cầu real-time gameplay
\end{itemize}
\introtext{Chiến lược bảo đảm dự án cTetris không chỉ đúng về mặt kỹ thuật mà còn mang lại trải nghiệm chơi hấp dẫn cho người dùng cuối.}
\label{LastPageTestStrategy}
\end{document}
